% 样本向量分解为常数方向的投影与正交残差（第4章）。示意图按示例数据 n=3 的长度绘制：
% x=(1,2,6)^T，样本均值 3，投影 3*1=(3,3,3)^T，长度 sqrt(27)≈5.196；
% 残差 r=(-2,-1,3)^T，长度 sqrt(14)≈3.742；||x||^2=41=27+14。
% 横轴代表一维常数方向 span{1}，纵轴代表 n-1 维残差子空间 {r: r_1+...+r_n=0}，二者正交。
\begin{tikzpicture}[line width=.6pt,font=\fontsize{8}{10}\selectfont,>=Stealth,x=8mm,y=8mm]
 \draw[FigureBlue,line width=.9pt] (-0.6,0)--(6.4,0) node[right,FigureBlue]{$\operatorname{span}\{\bm 1\}$};
 \draw[FigureGreen,line width=.9pt] (0,-0.4)--(0,4.6) node[anchor=west,FigureGreen]{残差子空间：$\sum_i r_i=0$，维数$n-1$};
 % 投影、残差与原向量
 \draw[->,FigureBlue,line width=1.1pt] (0,0)--(5.196,0);
 \node[FigureBlue,below] at (2.6,0) {$\bar x\bm 1$};
 \draw[->,FigureRed,line width=1.1pt,dashed] (5.196,0)--(5.196,3.742);
 \node[FigureRed,right] at (5.196,1.87) {$\bm r=\bm x-\bar x\bm 1$};
 \draw[->,FigureGreen,line width=1.1pt] (0,0)--(0,3.742);
 \node[FigureGreen,left] at (0,1.87) {$\bm r$};
 \draw[->,FigureInk,line width=1.1pt] (0,0)--(5.196,3.742);
 \node[FigureInk,above left] at (2.6,1.95) {$\bm x$};
 \draw[FigureInk,line width=.5pt] (4.896,0)--(4.896,0.3)--(5.196,0.3);
 \fill[FigureInk] (0,0) circle[radius=.5mm];
 \node[FigureInk,below left] at (0,0) {$\bm 0$};
 % 长度
 \node[FigureMuted,anchor=north] at (2.6,-0.45) {$\lVert\bar x\bm 1\rVert^2=n\bar x^2$};
 \node[FigureMuted,anchor=west] at (5.35,3.55) {$\lVert\bm r\rVert^2=\sum_i(x_i-\bar x)^2$};
\end{tikzpicture}
